\subsection{A balanced metric and a finite upper mass}
\label{sec:lower-mass-metric-geometry}

The metric used in the construction has additional differential properties. They distinguish this example from the Hopf threefold and identify a geometric obstruction that is independent of the scalar Monge--Amp\`ere equation. We retain the notation of \cref{eq:lower-mass-quotient-metric} and put $D=\mathrm{d}t\wedge\dc t$.

\begin{proposition}\label{prop:lower-mass-balanced-metric}
On the threefold of \cref{thm:vanishing-lower-mass}, the Hermitian form $\omega=\alpha+\beta+D$ satisfies
\[
 \mathrm{d}(\omega^2)=0,\qquad \ddc\omega=2\alpha\wedge\beta\geq0,
 \qquad \int_X\omega^3=24\pi.
\]
For every smooth real function $u$,
\begin{equation}\label{eq:lower-mass-vertical-energy}
 \int_X(\omega+\ddc u)^3
 =24\pi-6\int_X\mathrm{d}u\wedge\dc u\wedge\alpha\wedge\beta.
\end{equation}
Consequently $v_+(\omega)=24\pi$. In particular $X$ has the bounded mass property, while $v_-(h)=0$ for every Hermitian form $h$ on $X$.
\end{proposition}
\begin{proof}
Although $t$ is defined on the covering space, both $\mathrm{d}t$ and $\dc t$ descend. The Fubini--Study potentials give
\[
 B:=\ddc t=\beta-\alpha.
\]
The forms $\alpha,\beta$ are closed, and $\alpha^2=\beta^2=0$, since each is pulled back from a curve. As $D^2=0$, we obtain
\[
 \mathrm{d}D=-\mathrm{d}t\wedge B,\qquad
 \ddc D=-B^2=2\alpha\wedge\beta,
\]
\[
 \mathrm{d}(\omega^2)
 =-2(\alpha+\beta+D)\wedge\mathrm{d}t\wedge(\beta-\alpha)=0.
\]
In the last identity, the term involving $D$ has a repeated $\mathrm{d}t$, and
$(\alpha+\beta)\wedge(\beta-\alpha)=0$. Thus $\omega$ is balanced, hence Gauduchon, and the formula for $\ddc\omega$ follows.

On a fiber use $\zeta=a+\mathrm{i}b$ with both real periods $2\pi$. There $t=2a+$ a constant, so
\[
 D|_E=\frac1\pi\,\mathrm{d}a\wedge\mathrm{d}b,
 \qquad \int_ED=4\pi.
\]
The equality $\omega^3=6\alpha\wedge\beta\wedge D$, fiber integration and the integral-one normalization of each Fubini--Study form give $\int_X\omega^3=24\pi$.

Apply the Gauduchon threefold identity \eqref{eq:threefold-mass-dirichlet} with $\ddc\omega=2\alpha\wedge\beta$ to obtain \eqref{eq:lower-mass-vertical-energy}. Its last integral is nonnegative. Thus every smooth strictly $\omega$-plurisubharmonic mass is at most $24\pi$, and the constant potential attains this value. Bounded exact-background regularization, as in \cref{thm:gaud-volume-mass-extrema}, identifies the same supremum over bounded potentials. Metric independence in \cref{prop:bmp-metric} proves bounded mass on $X$. The assertion about the lower mass is \cref{thm:vanishing-lower-mass}.
\end{proof}

For an admissible smooth $u$, equality in the upper-mass bound holds precisely when $u$ is constant along every elliptic fiber. Indeed, $\alpha\wedge\beta$ is a positive volume form in the base directions, so the nonnegative integral in \eqref{eq:lower-mass-vertical-energy} vanishes exactly when the fiber derivative of $u$ vanishes. The fibers are connected, and smooth local sections then show that $u$ is the pullback of a smooth function on $\mathbf P^1\times\mathbf P^1$. For the sequence with masses tending to zero, that same integral instead tends to $4\pi$.

\begin{proposition}\label{prop:lower-mass-no-pluriclosed}
The threefold $X$ admits neither a pluriclosed Hermitian metric nor a closed K\"ahler current.
\end{proposition}
\begin{proof}
The nonzero smooth semipositive $(2,2)$-form $\alpha\wedge\beta$ is $\ddc$-exact by the preceding calculation. If a Hermitian metric $\eta$ satisfied $\ddc\eta=0$, Stokes would give
\[
 0<\int_X\eta\wedge\alpha\wedge\beta
   =\frac12\int_X\eta\wedge\ddc\omega
   =\frac12\int_X\omega\wedge\ddc\eta=0,
\]
a contradiction. The strict inequality holds because $\eta$ is positive in the vertical direction at every point.

If a closed positive current $T$ dominated $\varepsilon\omega$ for some $\varepsilon>0$, positivity and current Stokes would likewise give
\[
 0<\varepsilon\int_X\omega\wedge\alpha\wedge\beta
 \leq\int_XT\wedge\alpha\wedge\beta
 =\frac12\int_X\omega\wedge\ddc T=0.
\]
All pairings involve a current and smooth forms, so no singular product is being used. This excludes a closed K\"ahler current.
\end{proof}

The example is therefore outside the Fujiki class, by the characterization recalled in \cref{sec:non-kahler-geometric-problems}. It shows that a balanced metric does not guarantee positive lower mass, even together with bounded upper mass. In contrast, the Hopf threefold has positive lower mass and infinite upper mass. The two properties are logically independent. Their simultaneous failure on the products in \cref{cor:both-mass-properties-fail} is a third regime. In particular, bounded mass alone cannot replace the positive lower-mass hypothesis in a theorem asserting solvability for every non-pluripolar measure.
